% ------------------------------------------------------------
% Page 89
% ------------------------------------------------------------

A la suite de la Révocation de l’Edit de Nantes, en 1685, de nombreuses assemblées
clandestines sont tenues « \textit{aux appartenances du Mas de Conilhargues} ». Jean Roman,
prédicant très actif que l’on appelait souvent le Marchandou car il se déguisait en colporteur,
ou encore « \textit{le petit pacquetier} » ou « \textit{lou biquarel} », aurait convoqué plusieurs assemblées
au mas même.

\vspace{0.2cm}

L’une d’elles le \textbf{23 Novembre 1698}, est dénoncée par un nommé Poujol de Pourcarès. Les
dragons interviennent, mais sans succès, ne trouvant à la maison qu’Isabeau \textbf{Didès}, « \textit{perclue
et sourde} », qui assure n’avoir vu personne depuis plusieurs jours, « \textit{en raison de pluies
épouvantables qui inondent tous les chemins....} ». (Isabeau Dides était la belle mère de Nadal
Noël) Roux, fermier de Connilergues).
L’assemblée s’était déroulée, en fait, dans une grange de la ferme, Jean Roman ayant procédé
à plusieurs baptêmes et béni un mariage, devant environ 200 personnes.

\vspace{0.2cm}

Jacque Poujol, « \textit{le traître de Pourcarès} », fut exécuté par les camisards de la troupe de
Castanet en Octobre 1702, à cause de nombreuses dénonciations dont il fut l’auteur. On
retrouva son cadavre, sur le chemin, entre Pourcarès et Raffègues, le crâne fracassé par un
coup de pistolet, et accompagné d’un message écrit, mettant en garde d’éventuels
dénonciateurs, sur les risques qu’ils couraient...*

A cette époque le domaine appartenait à un sieur Pierre de Cros, seigneur de Campron, (sans
doute Campredon) ce qui expliquerait la présence d’un pigeonnier puisque seuls les nobles
avaient le privilège d’en posséder. En effet, selon les coutumes féodales il fallait être seigneur
d’un fief et exploitant d’un domaine pour avoir droit de « colombier », le fermier, lui, devant
subir les dégâts occasionnés par les pigeons. Le pigeonnier « sur pied » ou indépendant des
autres bâtiments était l’apanage des grands fiefs, ce qui était vraisemblablement le cas de
Connilergues.

\vspace{0.2cm}

Ce pigeonnier,* encore en bon état, la toiture exceptée, abrite quelques anciennes ruches à
cadres, délabrées. Il atteste ainsi, que les propriétaires de Connilergues disposaient de
ressources importantes pour l’époque. Ce qui expliquerait que le domaine ait été pillé à trois
reprises par les Camisards, lors de la guerre des Cévennes. Pierre de Cros était un « \textit{nouveau
converti} » comme M. de Thomassy de Meyrueis, qui s’était mis au service des troupes royales
« \textit{attirant ainsi sur lui, par un juste jugement} » les condamnations unanimes (?) et la haine de
quelques uns.

\vspace{0.2cm}

Nous disposons en effet, de « \textit{l’Etat des Nouveaux Convertis de la ville de Meyrueis et de
Gatuzières} » dressé en 1688 et décrivant leur attitude vis-à-vis de l’Eglise Romaine. Il est
question de Pierre De Cros, ou Du Cros, d’un de ses frères et d’un Hubac de Gatuzières sans
doute le meunier :

« \textit{Pierre De Cros ou Ducros, sieur de Campron, habitant faux bourg de l’église, 4 ou 5
enfants : Gabriel, Antoine, Marie, Madeleine, Marguerite. Son frère, sieur de Cambelongue
(?) et leur servante font très bien (leurs devoirs religieux de catholique). Il a beaucoup agi
pour empêcher les assemblées et en a interrompu une qui se tenait dans les bois près de
Cabrillac}. »

\vspace{0.3cm}

* Le pigeonnier a été remis en état début 2008
